%==============================================================================
%  CHAPTER 1 — GENERAL FRAMEWORK
%==============================================================================
\chapter{General Framework}
\label{chap:context}

\section{Introduction}
This chapter sets the frame of the project. It first presents the host
organization, then the framing of the project itself: its context, the problem
it answers, and the proposed solution. It closes with the working methodology,
the planning, and the tools used to manage the work.

% ---------------------------------------------------------------------------
\section{Presentation of the Host Organization}
\label{sec:host-org}

\subsection{Company Overview}
\textbf{INSOMEA} (INSOMEA Computer Solutions), a Beyon Solutions company, is a
Microsoft Cloud Gold Partner and one of the leading cloud IT solution providers
in the Middle East and Africa (MEA) region. Founded in 2016, it serves local,
regional, and international organizations and specializes in Microsoft~365,
Microsoft Azure, cybersecurity, and professional training.

The company operates from six offices across Tunisia, Bahrain, France, Algeria,
Morocco, and the Ivory Coast, with a team of more than fifty employees serving
over five hundred customers.

\begin{figure}[H]
  \centering
  
\includegraphics[height=1.9cm,keepaspectratio]{insomea}
  \caption{INSOMEA logo}
  \label{fig:insomea-logo}
\end{figure}

\subsection{Organizational Structure}
INSOMEA is organized around three delivery teams, one per activity pillar,
supported by a training center. This project was carried out inside the
\textbf{Cloud \& Infrastructure} team, as shown in
Figure~\ref{fig:org-structure}.

\begin{figure}[H]\centering
\begin{tikzpicture}[node distance=0.55cm and 0.45cm]
  \node[layernode,minimum width=5.2cm] (dir) {INSOMEA --- Management};
  \node[stepnode,below left=0.7cm and 2.6cm of dir.south,minimum width=3.1cm] (mw) {Modern\\Workplace};
  \node[stepnode,below=0.7cm of dir,minimum width=3.1cm] (sec) {Security};
  \node[stepnode,below right=0.7cm and 2.6cm of dir.south,minimum width=3.1cm,draw=diaggreen!70,fill=diaggreen!10] (cloud) {Cloud \&\\Infrastructure};
  \node[storenode,below=0.55cm of sec,minimum width=3.4cm] (tc) {INSOMEA\\Training Center};
  \node[stagenode,below=0.55cm of cloud,fill=diaggreen!12,draw=diaggreen!60,minimum width=3.1cm] (pfe) {This PFE\\(6 months)};
  \foreach \t in {mw,sec,cloud}{ \draw[flowarrow] (dir.south) -- (\t.north); }
  \draw[flowarrow,draw=diagamber!75] (sec.south) -- (tc.north);
  \draw[flowarrow,draw=diaggreen!70] (cloud.south) -- (pfe.north);
\end{tikzpicture}
\caption{Organizational structure and position of the project}
\label{fig:org-structure}
\end{figure}

\subsection{Activities and Areas of Expertise}
INSOMEA delivers IT solutions built on Microsoft technologies, organized around
three pillars:
\begin{itemize}[leftmargin=1.4em]
  \item \textbf{Modern Workplace}: collaboration and productivity with
        Microsoft~365 and Teams, including adoption and change management;
  \item \textbf{Security}: identity and access management, information
        protection and governance, threat protection;
  \item \textbf{Cloud \& Infrastructure}: Microsoft Azure, datacenter, DevOps,
        and application development and integration --- the team hosting this
        project.
\end{itemize}
The \textbf{INSOMEA Training Center} completes the offer with a broad catalog
of Microsoft certification courses delivered by Microsoft Certified Trainers.

\subsection{Key Partnerships}
INSOMEA's central partnership is with \textbf{Microsoft}: a Microsoft Cloud
Partner since 2017, the company has been named Microsoft \emph{Partner of the
Year} on five occasions, in Bahrain (2020) and in Tunisia (2021 and 2022)
among them, and holds more than fifteen Gold competencies together with eight
advanced specializations. As part of the \textbf{Beyon Solutions} group, it
also benefits from a regional network of technology companies across the MEA
region.

% ---------------------------------------------------------------------------
\section{Project Framing}
\label{sec:framing}

\subsection{Context and Problem Statement}
\label{sec:critique}
\gls{devops} unified software development and operations around automation and
\gls{cicd}. \gls{mlops} extends these principles to machine learning, where a
system is defined not by code alone but by \textbf{code}, \textbf{data}, and
\textbf{model} together: any of the three can change independently and break
the system. The resulting lifecycle forms a loop, shown in
Figure~\ref{fig:mlops-lifecycle}: upload, train, track, register, deploy,
serve, consume, monitor.

\begin{figure}[H]\centering
\begin{tikzpicture}
  \foreach \name [count=\n from 0] in {Upload,Train,Track,Register,Deploy,Serve,Consume,Monitor}{
    \node[stagenode] (s\n) at ({4.7*cos(90-\n*45)},{3.1*sin(90-\n*45)}) {\name};
  }
  \foreach \a/\b in {0/1,1/2,2/3,3/4,4/5,5/6,6/7,7/0}{
    \draw[flowarrow] (s\a) to[bend right=11] (s\b);
  }
  \node[align=center,font=\small\itshape,text=diagblue] at (0,0)
    {The end-to-end\\MLOps lifecycle};
\end{tikzpicture}
\caption{The end-to-end \gls{mlops} lifecycle automated by the platform}
\label{fig:mlops-lifecycle}
\end{figure}

In most teams, this lifecycle is still carried out by hand. A data scientist
trains a model in a notebook on a personal machine; results are recorded
informally, in the notebook output or a spreadsheet; the trained model is
shared as a file; an engineer wraps it in a web service through custom scripts;
and monitoring, if any, is limited to checking that the service responds.

This manual workflow suffers from well-known weaknesses:
\begin{itemize}[leftmargin=1.4em]
  \item \textbf{No reproducibility}: without a record of the exact code, data,
        and parameters, results cannot be reproduced;
  \item \textbf{Lost experiments}: informal tracking makes model comparison
        unreliable;
  \item \textbf{``Works on my machine''}: environment differences cause
        failures that are hard to diagnose;
  \item \textbf{Slow, fragile deployment}: hand-written scripts must be
        repeated for every model;
  \item \textbf{No isolation, no observability}: workloads share one machine
        and their health is invisible.
\end{itemize}

The problem of this project can therefore be summarized in one question:
\emph{how can a data scientist go from a notebook to a live, consumable
prediction service, without writing a single line of infrastructure code?}

\subsection{Proposed Solution}
\label{sec:solution}
The proposed solution is a self-hosted \gls{mlops} platform driven entirely
from a web browser. The user signs in, creates a \textbf{project}, uploads
code (\texttt{.py}, \texttt{.ipynb}, or \texttt{.zip}) and a dataset, and
triggers a \textbf{training pipeline} that runs in an isolated pod on a
\gls{k8s} cluster. Every run's metrics, parameters, and artifacts are captured
automatically.

From there, the user \textbf{registers} and promotes model versions
(Staging / Production / Archived), \textbf{deploys} a version as a live
prediction service in one click, tests it from the interface, and exposes it
to external applications through secured public \gls{api} keys. Dashboards,
serving telemetry, diagnostics, and an administration area complete the
platform. Figure~\ref{fig:solution-pipeline} summarizes this flow.

\begin{figure}[H]\centering
  \includegraphics[width=\textwidth,keepaspectratio]{solution_pipeline}
  \caption{Proposed solution: from upload to a live model with a public API}
  \label{fig:solution-pipeline}
\end{figure}

The platform runs on a real Kubernetes cluster and ships every change to a
public URL through an automated \gls{cicd} pipeline. It is assembled entirely
from open-source components; the study of existing platforms that motivates
this choice is the subject of Chapter~\ref{chap:study}.

% ---------------------------------------------------------------------------
\section{Methodology and Planning}
\label{sec:scrum}

\subsection{Scrum Methodology}
The platform is built incrementally from independent but connected
capabilities, so an \textbf{Agile} approach was preferred over a sequential
process, and \textbf{Scrum}~\cite{scrum_guide} was adopted. Scrum organizes
work into fixed-length \emph{sprints}, each producing a potentially shippable
increment, around three artifacts (product backlog, sprint backlog, increment)
and a set of events (planning, sprint, review, retrospective).
Figure~\ref{fig:scrum} illustrates the process.

\begin{figure}[H]
  \centering
  \includegraphics[width=0.9\textwidth,keepaspectratio]{scrum}
  \caption{The Scrum process: from product backlog to shippable increment}
  \label{fig:scrum}
\end{figure}

The Scrum roles were held as follows:
\begin{itemize}[leftmargin=1.4em]
  \item \textbf{\gls{po}}: Mr.\ Amine Gonji (company supervisor, INSOMEA),
        prioritizing the backlog and validating each increment;
  \item \textbf{\gls{sm}}: Mr.\ Ben Mardes Achref (academic supervisor),
        facilitating the process and removing impediments;
  \item \textbf{Development Team}: the author, responsible for analysis,
        design, implementation, and testing, with technical guidance from
        Mr.\ Benarous Abderrahmen.
\end{itemize}

\subsection{Sprint Organization}
The work was delivered over \textbf{four sprints in six months}, each ending
with a demonstrable, deployed increment reviewed with the supervisors
(Figure~\ref{fig:methodology-timeline}).

\begin{figure}[H]\centering
  \includegraphics[width=\textwidth,keepaspectratio]{sprints}
  \caption{The four sprints and their themes}
  \label{fig:methodology-timeline}
\end{figure}

\subsection{Product Backlog}
\label{sec:backlog}
The product backlog (Table~\ref{tab:backlog}) gathers all desired capabilities
as user stories, ordered by priority; story points estimate the relative
effort.

\begin{table}[H]\centering
\caption{Product backlog}\label{tab:backlog}
\renewcommand{\arraystretch}{1.15}
\small
\begin{tabularx}{\textwidth}{|c|X|c|c|}
\hline
\rowcolor{primary!15}
\textbf{\#} & \textbf{User story} & \textbf{Prio.} & \textbf{Pts}\\
\hline
1  & As a user, I want to sign up and sign in so that my work is private.        & High & 5\\ \hline
2  & As a user, I want to create and manage projects to organize my work.        & High & 5\\ \hline
3  & As a user, I want to upload code and a dataset through the browser.         & High & 5\\ \hline
4  & As a user, I want to trigger a training run on the cluster.                 & High & 8\\ \hline
5  & As a user, I want metrics and artifacts captured automatically.             & High & 5\\ \hline
6  & As a user, I want to watch the live logs of a running pipeline.             & Med  & 3\\ \hline
7  & As a user, I want to browse and compare my experiment runs.                 & Med  & 3\\ \hline
8  & As a user, I want to register and version trained models.                   & High & 5\\ \hline
9  & As a user, I want to promote a model between Staging/Production/Archived.   & High & 3\\ \hline
10 & As a user, I want to deploy a registered model as a live service.           & High & 8\\ \hline
11 & As a user, I want to send inputs and get predictions from the UI.           & High & 5\\ \hline
12 & As a user, I want a dashboard of cluster and per-pod metrics.               & Med  & 5\\ \hline
13 & As a user, I want public API keys so external apps can call my model.       & Med  & 5\\ \hline
14 & As a user, I want clear diagnostics when an uploaded script fails.          & Med  & 5\\ \hline
15 & As an administrator, I want to manage users and oversee the platform.       & Med  & 5\\ \hline
\end{tabularx}
\end{table}

The split into sprints is deliberate: each sprint delivers a complete vertical
slice that the previous one makes possible. Sprint~1 establishes identity and
ownership; Sprint~2 produces tracked training runs; Sprint~3 turns them into
registered, deployed models; Sprint~4 makes the whole lifecycle observable,
externally consumable, and administrable.

\subsection{Project Schedule}
The project ran over six months, from 12 January 2026, opening with a framing
and state-of-the-art phase, followed by the four sprints, with the report
written in parallel (Table~\ref{tab:schedule}).

\begin{table}[H]\centering
\renewcommand{\arraystretch}{1.2}
\caption{Project schedule}\label{tab:schedule}
\small
\begin{tabularx}{\textwidth}{|>{\bfseries}l|c|c|X|}
\hline
\rowcolor{primary!15}
\textbf{Phase} & \textbf{Period} & \textbf{Duration} & \textbf{Main objective}\\
\hline
Framing \& SOTA & 12 Jan -- 01 Feb & 3 weeks & Framing, state of the art, environment setup\\ \hline
Sprint 1 & 02 Feb -- 01 Mar & 4 weeks & Authentication and project management\\ \hline
Sprint 2 & 02 Mar -- 05 Apr & 5 weeks & Training pipeline and experiment tracking\\ \hline
Sprint 3 & 06 Apr -- 03 May & 4 weeks & Model registry, deployment, and prediction\\ \hline
Sprint 4 & 04 May -- 14 Jun & 6 weeks & Observability, public API, and administration\\ \hline
Report writing & 12 Jan -- 11 Jul & Continuous & Drafting and revising in parallel\\ \hline
\end{tabularx}
\end{table}

\subsection{Management and Communication Tools}
The project was managed with lightweight, traceable tools:
\begin{itemize}[leftmargin=1.4em]
  \item \textbf{Git and GitHub} for version control; every increment is a
        series of commits on the main branch, continuously deployed;
  \item \textbf{Visual Studio Code} as the development environment;
  \item \textbf{a work journal and weekly reports} (journal de bord, bilans)
        documenting progress, decisions, and blockers throughout the six
        months;
  \item \textbf{regular reviews} with the company and academic supervisors at
        the end of each sprint.
\end{itemize}

% ---------------------------------------------------------------------------
\section{Conclusion}
This chapter presented the host organization, framed the problem, one
question: from a notebook to a live prediction service with no infrastructure
code, introduced the proposed platform, and set the Scrum methodology, the
backlog, and the schedule that structure the rest of this report. The next
chapter analyzes the business objectives and studies the technical landscape.
